
\documentclass[11pt,a4paper]{article}
\usepackage[utf8]{inputenc}
\usepackage[T1]{fontenc}
\usepackage{lmodern}
\usepackage{amsmath,amssymb,mathtools}
\usepackage{geometry}
\usepackage{booktabs}
\usepackage{hyperref}
\usepackage{xcolor}
\usepackage{fancyhdr}
\geometry{margin=1.7cm}
\pagestyle{fancy}
\fancyhf{}
\lhead{NablaMath v0.3}
\rhead{Relatorio extremo}
\cfoot{\thepage}
\title{NablaMath v0.3: Relatorio matematico extremo}
\author{NablaMath}
\date{\today}
\begin{document}
\maketitle
\begin{abstract}
Este documento foi gerado por um modulo de relatorios matematicos em LaTeX. Ele combina expressao simbolica, derivadas passo a passo, gradiente, Hessiana, diagnostico numerico e manifesto de visualizacao.
\end{abstract}
\tableofcontents
\newpage
\section{Analise simbolica minuciosa}
\subsection{Expressao original}

\[
f=\left(x^{2} + y + -11\right)^{2} + \left(x + y^{2} + -7\right)^{2}
\]

A expressao e mantida como arvore simbolica, permitindo derivadas exatas, gradiente, Hessiana e avaliacao numerica.

\subsection{Gradiente}

\[
\nabla f=\begin{bmatrix}2\left(x^{2} + y + -11\right)2x + 2\left(x + y^{2} + -7\right)\\2\left(x^{2} + y + -11\right) + 2\left(x + y^{2} + -7\right)2y\end{bmatrix}
\]

\subsubsection{Derivada parcial em relacao a $x$}

\paragraph{Regra da soma} Derivamos cada parcela separadamente.

\[
\frac{d}{dx}\left(\left(x^{2} + y + -11\right)^{2} + \left(x + y^{2} + -7\right)^{2}\right)=\frac{d}{dx}\left(\left(x^{2} + y + -11\right)^{2}\right)+\frac{d}{dx}\left(\left(x + y^{2} + -7\right)^{2}\right)
\]

\paragraph{Regra da potencia com cadeia} Aqui u=x\textasciicircum{}\{2\} + y + -11 e n=2.

\[
\frac{d}{dx}u^n=n u^{n-1}u'
\]

\paragraph{Regra da soma} Derivamos cada parcela separadamente.

\[
\frac{d}{dx}\left(x^{2} + y + -11\right)=\frac{d}{dx}\left(x^{2} + y\right)+\frac{d}{dx}\left(-11\right)
\]

\paragraph{Regra da soma} Derivamos cada parcela separadamente.

\[
\frac{d}{dx}\left(x^{2} + y\right)=\frac{d}{dx}\left(x^{2}\right)+\frac{d}{dx}\left(y\right)
\]

\paragraph{Regra da potencia com cadeia} Aqui u=x e n=2.

\[
\frac{d}{dx}u^n=n u^{n-1}u'
\]

\paragraph{Regra do simbolo} A derivada de uma variavel em relacao a ela mesma e 1; em relacao a outra variavel e 0.

\[
\frac{d}{dx}x=1
\]

\paragraph{Regra do simbolo} A derivada de uma variavel em relacao a ela mesma e 1; em relacao a outra variavel e 0.

\[
\frac{d}{dx}y=0
\]

\paragraph{Regra da constante} Toda constante tem derivada zero.

\[
\frac{d}{dx}-11=0
\]

\paragraph{Regra da potencia com cadeia} Aqui u=x + y\textasciicircum{}\{2\} + -7 e n=2.

\[
\frac{d}{dx}u^n=n u^{n-1}u'
\]

\paragraph{Regra da soma} Derivamos cada parcela separadamente.

\[
\frac{d}{dx}\left(x + y^{2} + -7\right)=\frac{d}{dx}\left(x + y^{2}\right)+\frac{d}{dx}\left(-7\right)
\]

\paragraph{Regra da soma} Derivamos cada parcela separadamente.

\[
\frac{d}{dx}\left(x + y^{2}\right)=\frac{d}{dx}\left(x\right)+\frac{d}{dx}\left(y^{2}\right)
\]

\paragraph{Regra do simbolo} A derivada de uma variavel em relacao a ela mesma e 1; em relacao a outra variavel e 0.

\[
\frac{d}{dx}x=1
\]

\paragraph{Regra da potencia com cadeia} Aqui u=y e n=2.

\[
\frac{d}{dx}u^n=n u^{n-1}u'
\]

\paragraph{Regra do simbolo} A derivada de uma variavel em relacao a ela mesma e 1; em relacao a outra variavel e 0.

\[
\frac{d}{dx}y=0
\]

\paragraph{Regra da constante} Toda constante tem derivada zero.

\[
\frac{d}{dx}-7=0
\]

\paragraph{Resultado final da derivada} A expressao final e simplificada pelo motor simbolico da NablaMath.

\[
\frac{\partial}{\partial x}\left(\left(x^{2} + y + -11\right)^{2} + \left(x + y^{2} + -7\right)^{2}\right)=2\left(x^{2} + y + -11\right)2x + 2\left(x + y^{2} + -7\right)
\]

\subsubsection{Derivada parcial em relacao a $y$}

\paragraph{Regra da soma} Derivamos cada parcela separadamente.

\[
\frac{d}{dy}\left(\left(x^{2} + y + -11\right)^{2} + \left(x + y^{2} + -7\right)^{2}\right)=\frac{d}{dy}\left(\left(x^{2} + y + -11\right)^{2}\right)+\frac{d}{dy}\left(\left(x + y^{2} + -7\right)^{2}\right)
\]

\paragraph{Regra da potencia com cadeia} Aqui u=x\textasciicircum{}\{2\} + y + -11 e n=2.

\[
\frac{d}{dy}u^n=n u^{n-1}u'
\]

\paragraph{Regra da soma} Derivamos cada parcela separadamente.

\[
\frac{d}{dy}\left(x^{2} + y + -11\right)=\frac{d}{dy}\left(x^{2} + y\right)+\frac{d}{dy}\left(-11\right)
\]

\paragraph{Regra da soma} Derivamos cada parcela separadamente.

\[
\frac{d}{dy}\left(x^{2} + y\right)=\frac{d}{dy}\left(x^{2}\right)+\frac{d}{dy}\left(y\right)
\]

\paragraph{Regra da potencia com cadeia} Aqui u=x e n=2.

\[
\frac{d}{dy}u^n=n u^{n-1}u'
\]

\paragraph{Regra do simbolo} A derivada de uma variavel em relacao a ela mesma e 1; em relacao a outra variavel e 0.

\[
\frac{d}{dy}x=0
\]

\paragraph{Regra do simbolo} A derivada de uma variavel em relacao a ela mesma e 1; em relacao a outra variavel e 0.

\[
\frac{d}{dy}y=1
\]

\paragraph{Regra da constante} Toda constante tem derivada zero.

\[
\frac{d}{dy}-11=0
\]

\paragraph{Regra da potencia com cadeia} Aqui u=x + y\textasciicircum{}\{2\} + -7 e n=2.

\[
\frac{d}{dy}u^n=n u^{n-1}u'
\]

\paragraph{Regra da soma} Derivamos cada parcela separadamente.

\[
\frac{d}{dy}\left(x + y^{2} + -7\right)=\frac{d}{dy}\left(x + y^{2}\right)+\frac{d}{dy}\left(-7\right)
\]

\paragraph{Regra da soma} Derivamos cada parcela separadamente.

\[
\frac{d}{dy}\left(x + y^{2}\right)=\frac{d}{dy}\left(x\right)+\frac{d}{dy}\left(y^{2}\right)
\]

\paragraph{Regra do simbolo} A derivada de uma variavel em relacao a ela mesma e 1; em relacao a outra variavel e 0.

\[
\frac{d}{dy}x=0
\]

\paragraph{Regra da potencia com cadeia} Aqui u=y e n=2.

\[
\frac{d}{dy}u^n=n u^{n-1}u'
\]

\paragraph{Regra do simbolo} A derivada de uma variavel em relacao a ela mesma e 1; em relacao a outra variavel e 0.

\[
\frac{d}{dy}y=1
\]

\paragraph{Regra da constante} Toda constante tem derivada zero.

\[
\frac{d}{dy}-7=0
\]

\paragraph{Resultado final da derivada} A expressao final e simplificada pelo motor simbolico da NablaMath.

\[
\frac{\partial}{\partial y}\left(\left(x^{2} + y + -11\right)^{2} + \left(x + y^{2} + -7\right)^{2}\right)=2\left(x^{2} + y + -11\right) + 2\left(x + y^{2} + -7\right)2y
\]

\subsection{Hessiana}

\[
H_f=\begin{bmatrix}22x2x + 2\left(x^{2} + y + -11\right)2 + 2 & 22x + 22y\\22x + 22y & 2 + 22y2y + 2\left(x + y^{2} + -7\right)2\end{bmatrix}
\]

A Hessiana resume a curvatura local da funcao. Em otimizacao, ela ajuda a distinguir minimos, maximos e pontos de sela.


\section{Calculo numerico em malha}
Esta varredura numerica simples avalia a funcao em uma malha pequena. Ela nao substitui otimizacao global, mas oferece diagnostico inicial do dominio.
\begin{center}
\begin{tabular}{rrr}
\toprule
$x$ & $y$ & $f(x,y)$\\
\midrule
-6 & -6 & 890\\
-6 & -4 & 450\\
-6 & -2 & 610\\
-6 & 0 & 794\\
-6 & 2 & 810\\
-6 & 4 & 850\\
-6 & 6 & 1490\\
-4 & -6 & 626\\
-4 & -4 & 26\\
-4 & -2 & 58\\
-4 & 0 & 146\\
-4 & 2 & 98\\
-4 & 4 & 106\\
-4 & 6 & 746\\
-2 & -6 & 898\\
-2 & -4 & 170\\
-2 & -2 & 106\\
-2 & 0 & 130\\
-2 & 2 & 50\\
-2 & 4 & 58\\
-2 & 6 & 730\\
0 & -6 & 1130\\
0 & -4 & 306\\
0 & -2 & 178\\
0 & 0 & 170\\
0 & 2 & 90\\
0 & 4 & 130\\
0 & 6 & 866\\
2 & -6 & 1130\\
2 & -4 & 242\\
2 & -2 & 82\\
2 & 0 & 74\\
2 & 2 & 26\\
2 & 4 & 130\\
2 & 6 & 962\\
4 & -6 & 1090\\
4 & -4 & 170\\
4 & -2 & 10\\
4 & 0 & 34\\
4 & 2 & 50\\
4 & 4 & 250\\
4 & 6 & 1210\\
6 & -6 & 1586\\
6 & -4 & 666\\
6 & -2 & 538\\
6 & 0 & 626\\
6 & 2 & 738\\
6 & 4 & 1066\\
6 & 6 & 2186\\
\bottomrule
\end{tabular}
\end{center}
\[
\min_{grid} f \approx 10\quad \text{em}\quad (x,y)=(4,-2)
\]


\section{Metadados reprodutiveis}
\begin{verbatim}
{
  "generator": "NablaMath ExtremeLatexPDFReport",
  "created_at_unix": 1777159164.7620945,
  "expression": "(((((x^2) + y) + -11)^2) + (((x + (y^2)) + -7)^2))",
  "latex": "\\left(x^{2} + y + -11\\right)^{2} + \\left(x + y^{2} + -7\\right)^{2}",
  "variables": [
    "x",
    "y"
  ]
}
\end{verbatim}
\end{document}